\documentclass[11pt]{article}
\usepackage[a4paper,margin=23mm]{geometry}
\usepackage{amsmath,amssymb,array,booktabs,xcolor}
\usepackage[hidelinks]{hyperref}
\setlength{\parindent}{0pt}
\setlength{\parskip}{6pt}
\newcommand{\Q}{\mathbb Q}
\newcommand{\Z}{\mathbb Z}
\newcommand{\T}{\mathbb T}
\newcommand{\im}{\operatorname{im}}
\newcommand{\gr}{\operatorname{gr}}
\begin{document}
\begin{center}
{\Large A negative-cable counterexample to Problem 00000000614}\\[5pt]
{\large Actual diagrams and filtered rational Lee homology}\\[3pt]
lizaixi01\qquad 25 October 2026
\end{center}

\textbf{Claim and scope.}
The printed formula, with no restriction on the sign of $q$, is
\[
 s(C_{p,q}(K))=p\,s(K)+(p-1)(q-1).
\]
Take the crossingless unknot $U$ and $(p,q)=(2,-1)$. We compute both
invariants from their actual unreduced Lee complexes over $\Q$:
\[
 s(C_{2,-1}(U))=s(U)=0,\qquad
 2\cdot0+(2-1)(-1-1)=-2.
\]
Here $p>0$ and $\gcd(2,-1)=1$. This refutes only the unrestricted printed
formula. It does not refute a version with positive $q$ or another framing
convention.

\textbf{Published geometric boundary.}
We use the standard zero-framed satellite convention of
Lewark--Zibrowius, Definition 3.13: $t\in\Z$ and
$C_{2,2t+1}(U)=T_{2,2t+1}$; see printed pages 259--260.
Thus $t=-1$ is allowed. Identification of the product solid torus with the
zero-framed tubular neighbourhood of the unknot is this published geometric
identification. Lean does not reconstruct general three-manifold topology,
satellite invariance, Reidemeister invariance, or the Lee spectral sequence.
No such fact is an added Lean axiom. The particular torus curve, its entire
two-sheet braid, signs, finite incidence and resolutions are constructed
and checked in Lean.

\textbf{The actual negative diagram.}
On the genuine unit circle, use the embedded curve
\[
 \gamma:\T\longrightarrow\T\times\T,\qquad z\longmapsto(z^2,z^{-1}).
\]
It lies on the boundary of the product solid torus, is injective, and is
disjoint from the core $z\mapsto(z,0)$. Over longitude $e^{2\pi it}$,
$0\le t\le1$, its two sheets have meridians
$u_0(t)=e^{-\pi it}$ and $u_1(t)=-e^{-\pi it}$.
Lean proves these are all fibres, that they cover the entire curve, and
that their endpoints interchange.

Project to $(\Re u_s(t),-t)$, retaining $\Im u_s(t)$ as depth.
There is exactly one double point, $t=1/2$. Sheet 1 is above sheet 0.
The actual over/under tangents are $(\pi,-1)$ and $(-\pi,-1)$, whose
oriented crossing determinant is $-2\pi<0$. Therefore $n_+=0,n_-=1$.
The four ports are top left, top right, bottom left, bottom right,
numbered $0,1,2,3$. The strand involution is $(03)(12)$, the closure is
$(02)(13)$, and the clockwise rotation is $(0132)$. Their analytic
endpoint and tangent identifications are proved. The closed graph is
connected. In Rasmussen's Figure 1 conventions, the 0-resolution is
$(01)(23)$ and the 1-resolution is $(02)(13)$. Actual graph component
quotients give respectively one and two circles, hence a split saddle.

\newpage
\textbf{The rational Lee algebra and full complex.}
The native algebra is $A=\Q[\omega]/(\omega^2-1)$, implemented as
\texttt{QuadraticAlgebra}; $v_+=1$, $v_-=\omega$ form its native basis.
Actual multiplication and tensor comultiplication satisfy Rasmussen's
equations (3):
\[
 1\cdot1=\omega\cdot\omega=1,\quad
 1\cdot\omega=\omega\cdot1=\omega,\qquad
 \Delta(1)=1\otimes\omega+\omega\otimes1,\quad
 \Delta(\omega)=\omega\otimes\omega+1\otimes1.
\]
The actual circle-indexed tensor powers are linearly equivalent to $A$
and $A\otimes A$. Under these proved equivalences, the geometric saddle
is $\Delta$. The one-edge cube sign is $+1$. The full native complex,
indexed by every $i\in\Z$, is
\[
 C^{-1}=A\ \xrightarrow{\ \Delta\ }\ C^0=A\otimes A,
 \qquad C^i=0\quad(i\ne-1,0).
\]
It is a genuine \texttt{CochainComplex (ModuleCat Q) Z}, with $d^2=0$.
Injectivity of $\Delta$ is proved, so native homology is zero outside
degree 0. A canonical native homology isomorphism identifies
\[
 H^0(C)\simeq (A\otimes A)/\im\Delta.
\]

\textbf{Both normalization shifts are retained.}
For an enhanced state, Rasmussen uses
$\gr=|v|-n_-$ and $q=p+\gr+n_+-n_-$.
Thus the negative diagram's output has $q=p-1$, not $p$:
\[
\begin{array}{c|rr}
 C^{-1}\text{ basis}&1&\omega\\\hline
 q&-1&-3
\end{array}
\qquad
\begin{array}{c|rrrr}
 C^0\text{ basis}&00&01&10&11\\\hline
 q&1&-1&-1&-3 .
\end{array}
\]
These degrees follow from actual circle labels and actual crossing counts.
For every integer $j$, the cochain filtration consists of vectors whose
coefficients of basis vectors with $q<j$ vanish. Lean proves $\Delta$
preserves it. The homology filtration $F^jH$ is the image of its actual
filtered cycle representatives, transported through the canonical native
homology isomorphism.

\textbf{An explicit quotient calculation, not assigned homology values.}
Write a tensor vector as $x=(x_{00},x_{01},x_{10},x_{11})$.
Then
\[
 \Delta(a+b\omega)=(b,a,a,b),\qquad
 P(x)=(x_{00}-x_{11},x_{01}-x_{10}).
\]
The true equalities $\ker P=\im\Delta$ and surjectivity of $P$ are proved.
The section is $(a,b)\mapsto a(00)+b(01)$.
Consequently the native quotient has a proved coordinate isomorphism
to $\Q^2$; write $\overline P$ for this induced map. For every
$j\in\Z$, its actual filtered image is exactly
\[
 \overline P(F^jH)=\{(a,b):\ j>1\Rightarrow a=0,\quad
                          j>-1\Rightarrow b=0\}.
\]
The forward inclusion follows from all four actual tensor degrees;
the reverse inclusion uses the displayed genuine section. In particular,
the low-degree $11$ representative does not force its class to have
degree $-3$: the same class has the representative $-00$ of degree $1$.

\newpage
\textbf{Actual extrema and the crossingless companion.}
For a nonzero native homology class $h$, define its filtration degree by
\[
 d(h)=\sup\{j\in\Z:h\in F^jH\}.
\]
Lean proves each such set is nonempty and bounded above. In the actual
coordinates $(a,b)$ it is $(-\infty,-1]$ if $b\ne0$, and $(-\infty,1]$
if $b=0$ and $h\ne0$. Actual native classes with coordinates $(0,1)$
and $(1,0)$ attain both endpoints. Thus the set of nonzero-class degrees
is exactly $\{-1,1\}$; its minimum and maximum are genuine attained
extrema, not empty-set or unbounded default values.

The companion is the actual planar unit circle. Its projection is
injective and its native topological connected-component type has one
element. Its circle tensor space is proved equivalent to $A$. Its full
integer complex is $A$ in degree 0, zero elsewhere, with zero differential.
Native homology is $A$, with basis degrees $1,-1$, and all other native
homology groups are zero. The image-of-cycles filtration obeys the same
all-integer coordinate formula. Hence both diagrams have
\[
 s_{\min}=-1,\qquad s_{\max}=1,\qquad
 s=s_{\max}-1=s_{\min}+1=0
\]
in Definitions 3.1 and 3.4 of Rasmussen. This proves the stated
counterexample without assuming a precomputed unknot invariant.

\textbf{Formalization and reproducibility.}
The complete 969-line \texttt{lean/Main.lean} has 175 explicitly named
declarations, plus three anonymous instance commands. The dynamic
\texttt{Check.lean} enumerates every owned name, including generated and
anonymous instance names, and runs ordinary \texttt{\#print axioms} for
each. Only \texttt{propext}, \texttt{Classical.choice}, and
\texttt{Quot.sound} are allowed; there are no extra axioms, placeholders,
or \texttt{native\_decide} proofs.

The pins are Lean 4.33.0 and Mathlib commit
\texttt{db584cd6d46c92f209a44c0f1c829460d327499d}.
Run \texttt{lean --version}, then \texttt{bash lean/verify.sh} from the
appropriate project directory for build, warning-as-error declaration
audit, and the official fresh checker. The submission records the four
individual commands, source hashes and original resource observations.
Production elapsed times are diagnostic; they establish no controlled
research-efficiency improvement.

\textbf{Primary references actually used.}
\begin{enumerate}
\item J. Rasmussen, \emph{Khovanov homology and the slice genus},
arXiv:math/0402131v1 (9 February 2004), printed pp.3--4 and 6:
the resolutions, normalized cube, rational Lee maps and filtered
homology definitions.
\url{https://arxiv.org/abs/math/0402131}.
The un-oriented curl on p.15 is not used to assign this braid's sign.
\item L. Lewark and C. Zibrowius, \emph{Rasmussen invariants of Whitehead
doubles and other satellites}, J. reine angew. Math.816 (2024),241--296,
printed pp.259--260, Definition3.13 and its satellite convention.
\url{https://doi.org/10.1515/crelle-2024-0061}.
Only its geometric convention is used, not its characteristic-2 invariant.
\end{enumerate}
\end{document}
